% !TeX root = ../navier-stokes-manuscript.tex
\subsection{The two-dimensional derivative grid}\label{sub:SobolevRegularityInsideEveryCompletedRow}

Fix \(r\in\Nzero\).  The scalar energy of the response \(V_r\) at spatial
height \(s\) is
\[
  E_{r,s}:=\frac12\norm{\Lambda^sV_r}_2^2.
\]
The transfer paired with that response is
\[
  \mathcal P_{r,s}
  :=-
  \operatorname{Re}
  \pair{\Lambda^sV_r}
  {\Lambda^s\!\left[
    \sum_{a=0}^{r}\binom{r}{a}(V_a\cdot\nabla)V_{r-a}
    +\nabla Q_r
  \right]}_{L^2}.
\]
Pairing the \(r\)-th differentiated momentum equation at this height produces
\begin{equation}\label{eq:BodyEveryTemporalRowSpatialEnergy}
  E_{r,s}'=\mathcal P_{r,s}-2\nu E_{r,s+1}.
\end{equation}
Beginning at \(s=\tfrac12\) and repeating the induction proved in
\Cref{prop:CriticalHeightSpatialAscent} gives the completed row
\begin{equation}\label{eq:BodyEveryTemporalRowCompletedHeight}
  \partial_t^nE_{r,1/2}
  =(-2\nu)^nE_{r,n+1/2}
  +\sum_{j=0}^{n-1}
  (-2\nu)^j
  \partial_t^{\,n-1-j}\mathcal P_{r,j+1/2}.
\end{equation}
There is a spatial ladder inside \(u\), another inside \(u_t\), another
inside \(u_{tt}\), and one inside every finite response.  Moving down changes
the temporal response.  Moving across changes the spatial height.

\begin{center}
\captionsetup{hypcap=false}
\captionof{table}{The mixed derivative grid.  Each entry is a coordinate of \(u\); no entry is yet an endpoint estimate.}
\label{tab:MixedDerivativeGrid}
\renewcommand{\arraystretch}{1.35}
\begin{tabular}{@{}c c c c c c@{}}
\toprule
& \(H^0_x\) & \(H^1_x\) & \(H^2_x\) & \(H^3_x\) & \(\cdots\) \\
\midrule
\(V_0=u\) & \(\bullet\) & \(\bullet\) & \(\bullet\) & \(\bullet\) & \(\cdots\) \\
\(V_1=u_t\) & \(\bullet\) & \(\bullet\) & \(\bullet\) & \(\bullet\) & \(\cdots\) \\
\(V_2=u_{tt}\) & \(\bullet\) & \(\bullet\) & \(\bullet\) & \(\bullet\) & \(\cdots\) \\
\(V_3=u_{ttt}\) & \(\bullet\) & \(\bullet\) & \(\bullet\) & \(\bullet\) & \(\cdots\) \\
\(\vdots\) & \(\vdots\) & \(\vdots\) & \(\vdots\) & \(\vdots\) & \(\ddots\) \\
\bottomrule
\end{tabular}
\end{center}

The point of the grid is not to add its rows into an infinite norm.  Take one
finite request: \(q\) time derivatives and \(k\) spatial derivatives.  Choose
finite indices
\[
  r>q+\frac12,
  \qquad
  m>k+\frac32.
\]
The vector-valued time embedding followed by the three-dimensional spatial
embedding \parencite[Theorem~7.28]{Cerda2010} reads
\[
  H^r\bigl(I;H^m(\R^3)\bigr)
  \hookrightarrow
  C^q\bigl(\overline I;C^k(\R^3)\bigr).
\]
Every finite request has a finite coordinate above its thresholds.  Collecting
all of them on the fluid gives
\begin{equation}\label{eq:BodyFormalMixedIntersection}
  \mathscr V_\infty(I)
  :=\bigcap_{r=0}^{\infty}\bigcap_{m=0}^{\infty}
  H^r\bigl(I;H^m(\R^3)\bigr)
  \hookrightarrow C^\infty(\overline I\times\R^3).
\end{equation}
The constants may depend on the selected \((r,m)\).  Smoothness asks that
every finite coordinate exist; it does not ask for one constant uniform over
the entire grid.

The remaining issue is time.  A dot in the grid can be finite on every strict
subinterval and still diverge as the endpoint is approached.  To carry a
spatial coordinate to one time slice, we have to place it beside the temporal
response directly below it.  Those adjacent pairs determine the finite
rectangles.
